%!TEX program = xelatex
%!TEX encoding = UTF-8
\def\ano{2023}
\def\semestre{2}
\def\disciplina{Introdução à Álgebra Linear}
\def\turma{11}
\def\autor{José Antônio O. Freitas}
\def\instituto{MAT-UnB}

\documentclass{beamer}
\usetheme{Madrid}
\usecolortheme{beaver}
% \mode<presentation>
\usepackage[singlelinecheck=false]{caption}
\usepackage{textpos}
\usepackage{amssymb}
\usepackage{amsmath,amsfonts,amsthm,amstext}
\usepackage[justification=centering]{caption}
\usepackage[brazil]{babel}
% \usepackage[latin1]{inputenc}
\usepackage{graphicx}
\graphicspath{{../../figuras_geogebra/}{../../Pictures/}}
\usepackage{enumitem}
\usepackage{multicol}
\usepackage{answers}
\usepackage{qrcode}

\usepackage{pifont}
\newcommand{\cmark}{\ding{228}}%

% requires version 0.3 of the package
\usepackage[customcolors,beamer]{hf-tikz}

\tikzset{style green/.style={
        set fill color=green!50!lime!60,
        set border color=white,
    },
    style cyan/.style={
        set fill color=cyan!90!blue!60,
        set border color=white,
    },
    style orange/.style={
        set fill color=orange!80!red!60,
        set border color=white,
    },
    hor/.style={
        above left offset={-0.15,0.31},
        below right offset={0.15,-0.125},
        #1
    },
    ver/.style={
        above left offset={-0.1,0.3},
        below right offset={0.15,-0.15},
        #1
    }
}

\usepackage{tikz,ifthen,tikz-3dplot}
\usepackage{siunitx}
\usetikzlibrary{circuits.ee.IEC}
\usetikzlibrary{lindenmayersystems}
\usetikzlibrary[shadings]
\usepackage{pgfplots}
\pgfplotsset{compat=1.15}
\usepackage{mathrsfs}
\usetikzlibrary{arrows}
\usepackage{tcolorbox}
\newcounter{saveenumi}
\newcommand{\seti}{\setcounter{saveenumi}{\value{enumi}}}
\newcommand{\conti}{\setcounter{enumi}{\value{saveenumi}}}

\resetcounteronoverlays{saveenumi}


\renewcommand{\qedsymbol}{$\blacksquare$}

%\setbeamertemplate{itemize items}[ball]

%\setbeamertemplate{itemize items}[default]
%-------------amatrix
% Augmented matrix.  Usage (note the argument does not count the aug col):
% \begin{amatrix}{2}
    %   1  2  3 \\  4  5  6
    % \end{amatrix}
\newenvironment{amatrix}[1]{%
    \left[\begin{array}{@{}*{#1}{r}|r@{}}
    }{%
    \end{array}\right]
}

\newtheorem{definicao}{Defini\c{c}\~ao}[section]
\newtheorem{definicoes}{Defini\c{c}\~oes}[section]
\newtheorem{exemplo}{Exemplo}[section]
\newtheorem{exemplos}{Exemplos}[section]
\newtheorem{exercicio}{Exerc{\'\i}cio}
\newtheorem{observacao}{Observa{\c c}{\~a}o:}[section]
\newtheorem{observacoes}{Observa{\c c}{\~o}es:}[section]
\newtheorem*{solucao}{Solu{\c c}{\~a}o:}
\newtheorem{proposicao}{Proposi\c{c}\~ao}
\newtheorem{lema}{Lema}
\newtheorem{teorema}{Teorema}
\newtheorem{corolario}{Corol\'ario}
\newenvironment{prova}[1][Prova]{\noindent\textbf{#1:} }{\qedsymbol}%{\ \rule{0.5em}{0.5em}}
\def\ano{2023}
\def\semestre{2}
\def\turma{11}
\newcommand{\nsub}{\varsubsetneq}
\newcommand{\vaz}{\emptyset}
\newcommand{\dom}{{\rm dom\,}}
\newcommand{\im}{{\rm Im\,}}
\newcommand{\sub}{\subseteq}
\newcommand{\n}{\mathbb{N}}
\newcommand{\z}{\mathbb{Z}}
\newcommand{\rac}{\mathbb{Q}}
\newcommand{\real}{\mathbb{R}}
\newcommand{\complex}{\mathbb{C}}
\newcommand{\cp}[1]{\mathbb{#1}}
\newcommand{\ch}{\mbox{\textrm{car\,}}\nobreak}
\newcommand{\p}[1]{\mbox{\textrm{posto(#1)}}}
\newcommand{\vesp}[1]{\vspace{ #1  cm}}
\newcommand{\compcent}[1]{\vcenter{\hbox{$#1\circ$}}}
\newcommand{\comp}{\mathbin{\mathchoice
{\compcent\scriptstyle}{\compcent\scriptstyle}
{\compcent\scriptscriptstyle}{\compcent\scriptscriptstyle}}}
\newcommand\T{\rule{0pt}{2.6ex}}

\title{Diagonalização}
\author[\autor]{\autor}
\institute[\instituto]{\instituto}
\date{}

\begin{document}
    \begin{frame}
        \maketitle
    \end{frame}

    \logo{\includegraphics[scale=.1]{logo-MAT.png}\vspace*{8.5cm}}

    \begin{frame}
        Considere o operador linear $T \colon \real^3 \to \real^3$ dado por
        \begin{align*}
            T(x, y, z) = (x + 2y - z, -2x -3y + z, 2x + 2y - 2z).\pause
        \end{align*}

        Quando consideramos a base ordenada canônica \pause de $\real^3$, \pause $\mathcal{B} = \{e_1 = (1, 0, 0); e_2 = (0, 1, 0); e_3 = (0, 0, 1)\}$, \pause a matriz desse operador com respeito à essa base é\pause
        \[
            [T]_\mathcal{B} = \begin{bmatrix}
                \phantom{-}1 & \phantom{-}2 & -1\\
                -2 & -3 & \phantom{-}1\\
                \phantom{-}2 & \phantom{-}2 & \phantom{-}2
            \end{bmatrix}.
        \]
    \end{frame}

    \begin{frame}
        Agora, se tomamos a base ordenada \pause $\mathcal{B}_1 = \{w_1 = (1, 0, 2); w_2 = (0, 1, 2); w_3 = (1, -1, 1)\}$ \pause a matriz de $T$ com respeito à essa base é\pause
        \[
            [T]_{\mathcal{B}_1} = \begin{bmatrix}
            -1 & \phantom{-}0 & \phantom{-}0\\
            \phantom{-}0& -1 & \phantom{-}0\\
            \phantom{-}0& \phantom{-}0 & -2
        \end{bmatrix}.\pause
        \]

        Note que para a base $\mathcal{B}_1$ temos \pause
        \begin{align*}
            T(w_1) &= -1w_1\\
            T(w_2) &= -1w_2\\
            T(w_3) &= -2w_3\\
        \end{align*}
    \end{frame}

    \begin{frame}
        No caso geral, \pause seja $T : V \to V$ um operador linear \pause e suponha que exista uma base $\mathcal{B} \pause = \{v_1,\dots,v_n\}$ \pause de $V$ tal que\pause
        \begin{align}\label{formadiagonal}
            [T]_\mathcal{B} = \begin{bmatrix}
                \lambda_1 & 0 & 0 & \dots & 0\\
                0 & \lambda_2 & 0 & \dots & 0\\
                \vdots & & \ddots & & \vdots\\
                0 & 0 & 0 & \dots & \lambda_n
            \end{bmatrix}\pause
        \end{align}
        com $\lambda_i \in \cp{K}$ para $i = 1$, \dots, $n$.
    \end{frame}

    \begin{frame}
        Agora\pause
        \[
            [T(v_i)]_\mathcal{B}
            \begin{bmatrix}
            0\\
            0\\
            \vdots\\
            0\\
            \lambda_i\\
            0\\
            \vdots\\
            0
        \end{bmatrix}_\mathcal{B}\pause
        \]
        para $i = 1$, \dots, $n$. \pause Isto é,\pause
        \[
            T(v_1) = \lambda_1 v_1,\ \pause T(v_2) = \lambda_2 v_2,\pause \dots,\pause  T(v_n) = \lambda_n v_n.
        \]
    \end{frame}

    \begin{frame}
    \begin{definicao}
        Seja $T : V \to V$ um operador linear.\pause
        \begin{enumerate}[label={\roman*})]
            \item Um \textbf{autovalor} \pause de $T$ é um elemento $\lambda \in \cp{K}$ \pause tal que existe um vetor não nulo \pause $u \in V$ \pause com $T(u) = \lambda u$.\pause \index{Autovalor}
            \item Se $\lambda$ é um autovalor de $T$, \pause então todo vetor não nulo \pause $u \in V$ tal que\pause
            \[
                T(u) = \lambda u\pause
            \]
            é chamado de \textbf{autovetor} \pause de $T$ \textbf{associado} ao autovalor $\lambda$. \pause Denotaremos por $\aut_T(\lambda)$ \pause o subespaço gerado por todos os autovetores associados a $\lambda$. \pause Assim\pause
            \[
                \aut_T(\lambda) = \pause \{u \in V \pause \mid T(u) = \lambda u\}.\pause
            \]
            \item Suponha que $\dim V = n < \infty$. \pause Dizemos que $T$ é \textbf{diagonalizável} \pause se existir uma base $\mathcal{B}$ de $V$ \pause tal que $[T]_\mathcal{B}$ é diagonal, \pause isto é, tem a forma \eqref{formadiagonal}. \pause Tal fato equivale a dizer que existe uma base formada por autovetores.
        \end{enumerate}
    \end{definicao}
    \end{frame}

    \begin{frame}
        Temos aqui duas questões importantes:\pause
        \begin{enumerate}[label={\arabic*})]
            \item Dado um operador $T \colon V \to V$, \pause existem autovetores para $T$?\pause

            \item Se sim, \pause existe um número suficiente de autovetores \pause para forma uma base de $V$?
        \end{enumerate}
    \end{frame}

    \begin{frame}
        Aqui surgem duas questões:\pause
        \begin{enumerate}[label={\arabic*})]
            \item Como determinar se uma transformação linear $T \colon V \to V$ \pause possui autovetores?\pause

            \item A existência de autovetores de $T$ \pause depende da base escolhida para $V$?\pause
        \end{enumerate}
    \end{frame}

    \begin{frame}
        \begin{teorema}
            As seguintes afirmações são equivalentes:\pause
            \begin{enumerate}[label={\roman*})]
                \item o escalar $\lambda \in \cp{K}$ \pause é um autovalor do operador linear $T \colon V \to V$;\pause

                \item o operador $T - \lambda Id$ \pause não é injetor;\pause

                \item dada uma base ordenada $\mathcal{B}$ \pause qualquer de $V$ temos\pause
                \[
                    \det([T]_\mathcal{B} - \lambda[Id]_\mathcal{B}) = 0.
                \]
            \end{enumerate}
        \end{teorema}
    \end{frame}

    \begin{frame}
        \begin{teorema}
            Se $A \in M_n(\cp{K})$ \pause for uma matriz $n \times n$, \pause então $\lambda \in \cp{K}$ é um autovalor de $A$ \pause se, e somente se, $\lambda$ satisfaz a equação\pause
            \[
                \det(A - \lambda I_n) = 0.\pause
            \]
            Essa equação é chamada de \textbf{equação característica} de $A$.
        \end{teorema}
    \end{frame}

    \begin{frame}
        Seja $A \in M_n(\cp{K})$ onde $n \ge 1$. \pause Vamos definir o \textbf{determinante da matriz $A$}, \pause denotado por $\det(A)$, \pause de modo indutivo.\pause

        Se $n = 1$, \pause então a matriz $A \in \cp{M}_1(\cp{K})$ \pause é da forma
        \[
            A = (a_{11})\pause
        \]
        e neste caso definimos\pause
        \[
            \det(A) = a_{11} \in \cp{K}.
        \]
    \end{frame}

\begin{frame}
        Suponha que $n > 1$ \pause e que $\det(B)$ \pause esteja definido para todas as matrizes $B \in \cp{M}_p(\cp{K})$ \pause com $p < n$ \pause e seja $A \in \cp{M}_n(\cp{K})$. \pause Para cada $(i,j)$, \pause defina a matriz $A_{ij}$ \pause formada a partir de $A$ \pause retirando-se a sua $i$-ésima linha \pause e a sua $j$-ésima coluna. \pause É claro que $A \in \cp{M}_{n - 1}(\cp{K})$ \pause e portanto $\det(A_{ij})$ \pause está definido. \pause Defina então\pause
        \begin{align*}
            \det(A) &= \sum_{j = 1}^n(-1)^{i + j}a_{ij}\det(A_{ij}) \\ &= (-1)^{i + 1}a_{i1}\det(A_{i1}) + (-1)^{i + 2}a_{i2}\det(A_{i2}) \\ & + (-1)^{i + 3}a_{i3}\det(A_{i3}) + \cdots + (-1)^{i + n}a_{in}\det(A_{in}) \in \cp{K}
        \end{align*}
    \end{frame}

    \begin{frame}
        \begin{proposicao}
            Sejam $A$, $B \in \cp{M}_n(\cp{K})$ \pause e $\lambda \in \cp{K}$. \pause Temos:\pause
            \begin{enumerate}[label={\roman*})]
                \item $\det(AB) = \det(A) \det(B)$,\pause
                \item $\det(\lambda A) = \lambda^n \det(A)$,\pause
                \item Se $A$ é invertível, \pause então $\det(A^{-1}) = [(\det(A)]^{-1}$.
            \end{enumerate}
        \end{proposicao}
    \end{frame}

    \begin{frame}
        \begin{proposicao}
            Seja $A$ uma matriz $n \times n$ \pause com entradas num corpo $\cp{K}$.\pause
            \begin{enumerate}[label={\roman*})]
                \item Se $B$ é a matriz resultante da permutação de duas linhas de $A$, \pause então $\det (B) = -\det (A)$.\pause
                \item Se $B$ é a matriz resultante da multiplicação de uma linha de $A$ \pause por um escalar não nulo $\alpha \in \cp{K}$, \pause então $\det(B) = \alpha\det(A)$.\pause
                \item Se $B$ é a matriz resultante da soma da linha $i$ de $A$ \pause com um múltiplo não nulo \pause $\alpha \in \cp{K}$ da linha $j$ de $A$, \pause então $\det(B) = \det(A)$.
            \end{enumerate}
        \end{proposicao}
    \end{frame}

    \begin{frame}
        \begin{teorema}
            Uma matriz $A \in \cp{M}_n(\cp{K})$ \pause é invertível se, \pause e somente se, $\det(A) \ne 0$.
        \end{teorema}
    \end{frame}

    \begin{frame}
        Dado $V$ um espaço vetorial sobre $\cp{K}$, \pause com $\dim_\cp{K} V = n < \infty$ \pause e um operador linear $T \colon V \to V$, \pause queremos tentar encontrar, se existir, \pause uma base ordenada $\mathcal{D}$ de $V$ \pause formada por autovetores. \pause Para isso devemos proceder da seguinte maneira:\pause
        \begin{enumerate}[label={\arabic*})]
            \item Se não foi dada, \pause encontre a matriz de $T$ com relação a alguma base ordenada $\mathcal{B}$ de $V$. \pause Pode-se usar a base canônica $\mathcal{B}$ \pause de $V$ para facilitar o cálculo de $[T]_\mathcal{B}$.
            \seti
        \end{enumerate}
    \end{frame}

    \begin{frame}
        \begin{enumerate}[label={\arabic*})]
            \conti
            \item Uma vez encontrada $[T]_\mathcal{B}$ calcule\pause
            \[
                [T]_\mathcal{B} - \lambda I_n = [T]_\mathcal{B} - \begin{bmatrix}\lambda & 0 & 0 & \cdots & 0\\0 & \lambda & 0 & \cdots & 0\\0 & 0 & \lambda & \cdots & 0\\\vdots & \vdots & \vdots & \cdots & \vdots\\0 & 0 & 0 & \cdots & \lambda\end{bmatrix}.\pause
            \]

            \item\label{passo4diag} Determine todos os valores, \pause se existem, de $\lambda \in \cp{K}$ \pause tais que\pause
            \[
                \det([T]_\mathcal{B} - \lambda I_n) = 0.
            \]
            \seti
        \end{enumerate}
    \end{frame}

    \begin{frame}
        \begin{enumerate}[label={\arabic*})]
            \conti
            \item Se não existir $\lambda \in \cp{K}$ \pause no passo anterior, \pause então $T$ não é diagonalizável.\pause

            \item Para cada $\lambda_1$, \pause $\lambda_2$, \pause \dots, $\lambda_r$ \pause encontrando em \ref{passo4diag} \pause resolva o sistema linear homogêneo\pause
            \[
                ([T]_\mathcal{B} - \lambda_jI_n)X = 0.\pause
            \]
            O subespaço $Aut_T(\lambda_j)$ \pause será dada pelas soluções desse sistema.
            \seti
        \end{enumerate}
    \end{frame}

    \begin{frame}
        \begin{enumerate}[label={\arabic*})]
        \conti
            \item Encontre uma base para cada $Aut_T(\lambda_j)$, \pause com $j = 1$, 2, \dots, $r$.\pause

            \item Se o conjunto $\mathcal{D}$ \pause resultante a união de todas as bases encontradas no passo anterior \pause possuir $\dim_\cp{K}V$ elementos, \pause então $\mathcal{D}$ será uma base de $V$ \pause tal que a matriz de $T$ com relação à essa base será diagonal. \pause Se $\mathcal{D}$ possuir menos elementos que $\dim_\cp{K}V$ \pause então $T$ não é diagonalizável.
        \end{enumerate}
    \end{frame}

    \begin{frame}
        \begin{teorema}
            Seja $T : V \to V$ um operador linear onde $V$ é um $\cp{K}$-espaço vetorial de dimensão finita e sejam $\lambda_1$, \dots, $\lambda_r$, $r \ge 1$, autovalores de $T$, dois a dois distintos.
            \begin{enumerate}[label={\roman*})]
                \item\label{autovetorLI} Se $u_1 + \cdots + u_r = 0_V$ com $u_i \in \aut_T(\lambda_i)$; $i = 1$, \dots, $r$; então $u_i = 0_V$ para todo $i$.
                \item Para cada $i = 1$, \dots, $r$ seja $\mathcal{B}_i$ um conjunto linearmente independente contido em $\aut_T(\lambda_i)$. Então $\mathcal{B}_1 \cup \cdots \cup \mathcal{B}_r$ é L.I. em $V$.
            \end{enumerate}
        \end{teorema}
    \end{frame}

    \begin{frame}
        \begin{corolario}
            Seja $T : V \to V$ um operador linear, onde $V$ é um $\cp{K}$-espaço vetorial de dimensão finita. Se $\lambda_1$, \dots, $\lambda_r$ são todos os autovalores de $T$, então $T$ é diagonalizável se, e somente se,
            \[
                \dim_\cp{K} V = \sum_{i = i}^r\dim_\cp{K}\aut_T(\lambda_i).
            \]
        \end{corolario}
    \end{frame}
\end{document}
